\section{Unit Propagation}

Unit propagation is the next step in implementing DPLL, and it introduces the
idea of implication. Imagine a clause with two variables: $A$ and $Z$,
and we know that $A = False$. What can we say about $Z$? Well, if our equation is satisfiable, 
then the clause must be satisfiable, therefore $Z = True$ because if it didn't, 
that would be a contradiction.

So instead of assigning variables $B$, $C$, \ldots, $Y$, then assigning $Z$,
why don't we just do it now, and use that knowledge in other clauses
containing $Z$ or $\neg Z$?

That is essentially the idea behind unit propagation: the decisions we make
imply other assignments, and so we should force them now instead of waiting
until later.

However, this introduces a light hiccup into our implementation. You see, 
our initial invariant about assigning variables is that they are
assigned in numerical order instead of (possibly) out of order.

``Not a problem! We'll just make the assignment array have space for every
variable so we can index it anywhere,'' I hear you think. Unfortunately that interferes with
our base case because it relies on the size
of the assignment array. If it is initially at a fixed size of
\textit{numVars}, then our base case will always fire because our assignment
array always has a size of \textit{numVars}. So we would either need to
change our base case or change how we store learned variables.

What about a learned set? One that holds whether or not we've learned a
variable (positive or negative), and then as we traverse down the call stack,
we just check if a variable is already learned, and if so, we give it its learned
value?

That would 100\% work. In fact, the first time I implemented DPLL,
that is exactly what I used. However, it has the disadvantage of being annoying 
to maintain. Instead, why don't we create a new data type called something like
\textit{var} that represents the state of each variable: \textit{True},
\textit{False}, and \textit{Unassigned}.

Then we can just maintain a simple counter of how many are not
\textit{Unassigned}, and then that keeps our base case nice and neat.
Yes dear reader, we are going to do that, however, if you'd like to use the
learned set, reading about how we do it without a learned set will help you
and offer clues on how to implement it your way.

\subsection{Changes To Our Implementation}

\begin{lstlisting}
    public interface:
        SATSolver(string input)
        bool solveCNF()

    private interface:
        enum var
        {
            False
            True
            Unassigned
        }

        void parse(string equation)
        bool solve(var[] assignment, int level)
        bool evaluate(var[] assignment)
        void propagate()

        int numVars
        int numClauses
        int numAssigned

        int[][] clauses

        optional var[] vars
\end{lstlisting}

Here we just add in a new enum for the states that a variable can be, and
adjust the methods using \textit{bool}s to use the new enum type instead.

\subsection{``Textbook" Unit Propagation}

Before we start changing our algorithm, let's implement the
\textit{propagate()} method.

As stated before, the idea is that given a clause with either only one variable
or only one \textit{Unassigned} variable (and the rest evaluate to
\textit{False}), we can immediately say that \textit{Unassigned} variable will
be whatever value it needs to be to evaluate to \textit{True} for that clause.

In the ``textbook" way of implementation,
\footnote{\url{https://en.wikipedia.org/wiki/Unit_propagation}} 
once a clause is satisfied because of a literal, you remove every other
clause that is satisfied by that literal. And for every other clause containing 
the negation of the literal, you remove that negation entirely. This is done to shrink the 
search space.

We will implement this algorithm with as much malicious compliance as you can 
carry at an office job without getting fired:

\begin{algorithm}[H]
\caption{Main Solving Method}\label{alg:9}
\begin{algorithmic}[1]

\Procedure{solveCNF}

\State assignment[1..numVars] \Comment{assignment contains numVars variables}

\For{each $var$ in $assignment$}
    \State $var \gets Unassigned$
\EndFor

$\Call{Propagate}{assignment}$

\State return $\Call{solve}{assignment, 1}$

\EndProcedure
\end{algorithmic}
\end{algorithm}

We want to propagate before we start solving in case there are clauses with 
only one variable (they're already unit). One thing to note about this pseudocode 
implementation that differs in \textbf{Algorithm 8} is that here, $assignment$ 
already contains $numVars$ variables, whereas in \textbf{Algorithm 8}, $assignment$
contained up to $n$ variables. They both still use 1-based indexing.

\begin{algorithm}[H]
\caption{Recursive Solving Method}\label{alg:10}
\begin{algorithmic}[1]

\Procedure{solve}{assignment[1..n], level}

\If{level = numVars}
    \State return $\Call{evaluate}{assignment}$
\EndIf

\State $backupAssignment[1..numVars] \gets assignment$
\State $backupClauses[1..numClauses][1..n] \gets clauses$

\If{assignment[level] = Unassigned}
    \State $assignment[level] \gets true$
    $\Call{propagate}{assignment}$
\EndIf

\If{$\Call{solve}{assignment, level + 1} = true$}
    \State return $true$
\EndIf

\State $clauses \gets backupClauses$
\State $assignment \gets backupAssignment$

\If{assignment[level] = Unassigned}
    \State $assignment[level] \gets false$
    $\Call{propagate}{assignment}$
\EndIf

\If{$\Call{solve}{assignment, level + 1} = true$}
    \State return $true$
\EndIf

\State{return $false$}

\EndProcedure
\end{algorithmic}
\end{algorithm}

Before we get onto the code for unit propagation, let's understand this code first. 
Because unit propagation no longer guarantees that variables are assigned in order, 
we need to pass in a ``level'' as a parameter so we know which variable we need to 
assign. We also now need to check that propagation didn't already assign this variable. 
If it's currently Unassigned, we need to assign it then propagate, but if it's been 
assigned before we assign it, then we know that it's the result of propagation, so we 
can continue onto the next level. Other than that, it's practically the same as our recursive 
backtracking code, except now we need to have a ``backup'' of our clauses and our assignment. 

Why? Because \textit{propagate()} \textbf{mutates} the clauses and the assignment in a way 
that breaks the assumption that our literals are assigned in ``numeric'' order. Since it 
returns nothing, we don't have a way to undo that easily, so instead, for readability, 
we just have a backup copy and whenever propagation leads to an unsatisfying assignment, 
we return false and undo whatever changes we made by copying back into our solver's fields.
And the reason why we don't need to undo the propagation before returning False is because 
when we move back up the call stack by one level, the method continues by undoing the effects for us.
Now let's look at how we perform textbook unit propagation:

\begin{algorithm}[H]
\caption{Textbook Unit Propagation}\label{alg:11}
\begin{algorithmic}[1]

\Procedure{propagate}{assignment[1..n]}

\While{$true$}
    \State unitClauses \Comment{A set of clauses}
    \State unitVars \Comment{A set of variables}

    \For{each $Clause$ in $Clauses$}
        \State $isUnit \gets true$
        \State $uIdx \gets -1$
        \State $numU \gets 0$

        \For{each $var$ in $Clause$}
            \State $idx \gets \Call{abs}{var}$

            \If{assignment[$idx$] = Unassigned}
                \State $numU \gets numU + 1$
                \State $uIdx \gets \Call{indexOf}{Clause, var}$
            \ElsIf{$\Call{evalsFalse}{var} = false$}
                \State $isUnit \gets false$
            \EndIf
        \EndFor

        \If{isUnit $\wedge$ numU = 1}
            \State \Call{insert}{$unitClauses$, $Clause$}
            \State \Call{insert}{$unitVars$, $uIdx$}

            \State $var \gets clause[uIdx]$
            \State $idx \gets \Call{abs}{var}$

            \If{$var > 0$}
                \State $assignment[idx] \gets true$
            \Else 
                \State $assignment[idx] \gets false$
            \EndIf
        \EndIf
    \EndFor

    \For{each $Clause$ in $Clauses$}
        \For{each $var$ in $Clause$}
            \If{$\Call{contains}{unitVars, -var}$}
                \State $\Call{remove{Clause, -var}}$
            \EndIf

            \If{$\Call{contains}{unitVars, var} \wedge \Call{contains}{unitClause, Clause}$}
                \State $\Call{remove{Clauses, Clause}}$
            \EndIf
        \EndFor
    \EndFor

    \If{\Call{Empty}{unitClauses} $\wedge$ \Call{Empty}{unitVars}}
        return
    \EndIf
\EndWhile

\EndProcedure
\end{algorithmic}
\end{algorithm}

Okay, now let's understand what's happening here. We immediately go into an 
infinite loop (nice), and we start looking at each clause for two things:
\begin{enumerate}
    \item The clause has only one Unassigned variable 
    \item Every other literal in the clause is False (remember, literals are evaluated variables)
\end{enumerate}

If both things are true, then the Unassigned variable will be assigned whatever value 
is required to make it evaluate to True within that clause, and we know we have found 
a unit variable and a unit clause. After we have found all of our unit variables and unit 
clauses, we then rescan every other clause and delete the each unit variable's negative,
(if we learn that -5 is unit, then we remove +5), and if the clause contains a unit 
variable, we just delete that clause entirely from the database. Once we stop finding 
new unit variables and new unit clauses, we stop entirely. You may wonder (I reread my 
code and thought I found a bug), ``how do we know that we always terminate?'' That's a 
great question. You see, this confusion may come from the fact that we rescan the database 
every iteration of the \textit{while()} loop searching for unit clauses. What is stopping 
us from re-finding the same clause as unit upon a new iteration? The answer is much simpler 
than you'd expect: since we assign our unit variables, the clauses upon the next iteration 
no longer foloow our criteria for being unit propagate-able because there is now a literal 
in that clause that is True. Therefore, the number of unit variables and clauses monotonically
decrease between iterations (a fancy way of saying the number can only go down).

Now, you may be wondering, ``what's the performance of our solver now?" Let me answer that 
as plainly as possible: worse than brute force. How? Easy! We are doing SO MANY things that 
hurt the computer. We are allocating TONS of memory in order to backtrack, we are allocating 
even more memory to propagate (hash sets aren't cheap), and hashing, while O(1), is not 
a cheap operation to do. So yeah, I can believe that it's worse than brute force because 
at least the brute force algorithm was easy on the CPU. Let's learn how this is 
\textit{usually}
\footnote{I wanted to say ``actually implemented'', however, that's a big generalization, and 
again, since I'm not the expert here, I don't think I have the authority to say that. The next 
section is how \textit{I} implemented it, and it roughly follows the idea on how SAT solvers 
implement this.}
implemented!
